% New review increment. No supplied source is modified or included automatically.
% Label prefix: inc:AN06:v1:
\documentclass[11pt]{article}
\usepackage[T1]{fontenc}
\usepackage[utf8]{inputenc}
\usepackage{lmodern}
\usepackage[margin=0.9in]{geometry}
\usepackage{amsmath,amssymb,amsthm,mathtools,booktabs,array,microtype}
\usepackage{xurl}
\usepackage[hidelinks]{hyperref}
\usepackage{fancyhdr}
\pagestyle{fancy}\fancyhf{}
\fancyhead[L]{\small AN06: learned limiting-orbit witnesses}
\fancyhead[R]{\small Review increment v1}
\fancyfoot[C]{\thepage}
\setlength{\headheight}{14pt}
\setlength{\emergencystretch}{3em}
\allowdisplaybreaks[2]
\numberwithin{equation}{section}
\newtheorem{theorem}{Theorem}[section]
\newtheorem{lemma}[theorem]{Lemma}
\newtheorem{proposition}[theorem]{Proposition}
\newtheorem{corollary}[theorem]{Corollary}
\theoremstyle{remark}\newtheorem{remark}[theorem]{Remark}
\newcommand{\E}{\mathbb E}
\newcommand{\R}{\mathbb R}
\newcommand{\dd}{\,d}
\newcommand{\1}{\mathbf 1}
\newcommand{\norm}[1]{\left\lVert#1\right\rVert}
\title{Learned witnesses on the selected coherent orbit\\
\large Half-learning descent, common competition windows, and population robustness}
\author{Research proof increment for independent review}
\date{Version 1 --- October 2, 2026}
\begin{document}
\maketitle
\begin{abstract}
For the selected coherent weak-growth orbit of the defined leading population
system, we prove a coupled invariant region through weak half-mass, with
$y'(0)\le-5/496$ and $\mathcal E_\zeta'(0)\le-1/250$ uniformly for
$\rho\in[13/20,7/10]$ and $\zeta\in[4,5]$.
Consequently a competition crossover can be chosen strictly after half-mass.
We construct common positive-width time windows and a local parameter/probe
rectangle by explicitly defined analytic constants, without assuming a simple
or unique crossing. We also give a quantitative first-moment distance bound
for the leading population rates, including nonunit strong mass and probe-gate
mismatches, and a conditional finite-window propagation estimate.
A separate static Gaussian estimate converts the reached mass bounds into
exact response and cluster-loss bounds for a fully charted finite-noise state.
None of these statements proves isotropic entry, nonlinear long-delay passage,
or finite-noise trajectory/rate transfer. The initialized Theorems A and B
remain open. No numerical trajectory is a proof input.
\end{abstract}

\section{Proof contract and notation}\label{inc:AN06:v1:contract}
\paragraph{Dependencies.}
The working base is \path{THEOREM_A_ANALYTICAL_PROGRESS.tex}, version 1.0.
We use its \texttt{pa:branch} (selected coherent branch),
\texttt{pa:weakroot} (real resident weak root),
\texttt{pa:atomiclag} and \texttt{pa:atomicrates} (lag and source identities),
and \texttt{pa:limitreversal} (global continuation, $v>0$, $v>y$, $v>x$,
and eventual hit of every finite probe edge greater than $a_\rho$).
The leading population equations are the defined system in
\texttt{pa:leading}/\texttt{pa:masslaws}, not an assumed theorem of
finite-noise convergence. Relevant equations are restated below.
The original initialization/output identity used only in the static learning
comparison is $f_0(X)=S_0X/4$.

\paragraph{Scope.}
Sections 2--7 concern the \emph{defined leading system}.
Section 8 is a conditional, static result for the exact Gaussian model.
No source, AN02 increment, or consolidated checkpoint is overwritten.
The interval $[4,5]$ of scaled probe offsets is an analytical witness choice;
it does not include the reported canonical scaled offset near $1.745$.
The common-window parameter interval constructed below is local, not the
entire half-learning rectangle.

Set $\lambda=\rho^2$, and let $\phi,\Phi$ be the standard Gaussian density
and distribution function. Define
\[
 K(r)=\phi(r)+r\Phi(r),\qquad g(r)=K(r)-r>0.
\]
For the selected coherent orbit write $Z_\rho=(n,x,y,u,v)$, where
\begin{align}
 n'&=\lambda n(1-n),\qquad n(s)=(1+e^{-\lambda s})^{-1},
 \label{inc:AN06:v1:orbit}\\
 x'&=\tfrac\lambda2[nv-x+(1-n)y]\Phi(\rho x),\nonumber\\
 y'&=\tfrac12[-y-nK(u)+\rho(1-n)K(\rho x)],\nonumber\\
 u'&=\tfrac12[(1-n)(\phi(u)-\lambda u)-(nu+y)\Phi(u)],\nonumber\\
 v'&=\tfrac12[\rho K(\rho x)-\lambda v].\nonumber
\end{align}
The branch approaches
\[
 (0,a,a,u_*,a/\lambda)\quad(s\to-\infty),\qquad
 a=\rho K(\rho a),\quad \phi(u_*)-a\Phi(u_*)-\lambda u_*=0.
\]
Time $s=0$ means weak half-\emph{mass}; exact finite-noise learning is
addressed separately in Section 8.
For $x<\zeta$, put
\begin{align}
 \mathcal E_\zeta&=-\tfrac12\zeta^2+\tfrac12x^2+y(\zeta-x),
 &\mathcal E_\zeta'&=(x-y)x'+(\zeta-x)y',
 \label{inc:AN06:v1:error}\\
 L&=y+nK(u)>0,
 &d_1&=-\tfrac12(\zeta-x)L,\nonumber\\
 d_2&=(x-y)x'+\tfrac\rho2(\zeta-x)(1-n)K(\rho x),
 &\mathcal E_\zeta'&=d_1+d_2.\nonumber
\end{align}
Throughout, a prime is differentiation in normalized/shifted time, not
physical time. The latter would multiply an original finite-noise rate by
$\mu_1^2$; no finite-noise rate approximation is assumed here.

\section{A coupled region through weak half-mass}
\label{inc:AN06:v1:halfregion}
Fix
\begin{equation}\label{inc:AN06:v1:rhointerval}
 \mathcal K=[13/20,7/10],\qquad
 169/400\le\lambda\le49/100.
\end{equation}
Use the following elementary Gaussian bounds:
\begin{equation}\label{inc:AN06:v1:gaussianbounds}
 \frac38<\phi(0)<\frac25,\qquad
 \frac15<\phi(1)<\frac14,\qquad
 K(r)\le\phi(0)+\frac r2+\frac{\phi(0)r^2}{2}\quad(r\ge0).
\end{equation}
For completeness, the first two follow from the elementary bounds
$25/8<\pi<22/7$ and $8/3<e<3$; the last follows by integrating
$\Phi(r)\le1/2+\phi(0)r$. Also $K(r)\le r_++\phi(0)$ for all real $r$.

\begin{lemma}[Resident bounds]\label{inc:AN06:v1:residentbounds}
For $\rho\in\mathcal K$, $0<a<\phi(0)<2/5$ and $0<u_*<1$.
\end{lemma}
\begin{proof}
The function $z-\rho K(\rho z)$ is strictly increasing. At $z=\phi(0)$,
\[
 \frac{\rho K(\rho\phi(0))}{\phi(0)}
 \le\rho+\frac{\rho^2}{2}+\frac{\rho^3\phi(0)^2}{2}
 \le\frac7{10}+\frac{49}{200}+\frac{343}{12500}<1.
\]
Its value at zero is negative, proving the first claim. The weak-root
function is positive at zero since $\phi(0)-a/2>0$, and negative at one
since $\phi(1)<1/4<\lambda$. Uniqueness of the real weak root proves the
second claim.
\end{proof}

Define
\begin{equation}\label{inc:AN06:v1:TZ}
 A=\rho K(\rho x),\quad P=\Phi(\rho x),\quad Q=\Phi(u),\quad
 z=x-y,\quad T=L-(1-n)A=-2y'.
\end{equation}
The letter $T$ here is a scalar output-descent observable, not an entry time.

\begin{theorem}[Half-learning region and output descent]
\label{inc:AN06:v1:regiontheorem}
On the selected branch, for every $\rho\in\mathcal K$ and every $s\le0$,
\begin{equation}\label{inc:AN06:v1:region}
 -n<y<a<\frac25,\quad y<x<\frac34,\quad
 0<u<1,\quad 0<v<\frac{11}{10},\quad T>0.
\end{equation}
Moreover,
\begin{equation}\label{inc:AN06:v1:Tlower}
 T(s)\ge\frac5{62}n(s)(1-n(s)),\qquad
 y'(0)\le-\frac5{496}.
\end{equation}
\end{theorem}
\begin{proof}
\emph{Entry at the resident end.}
The checkpoint's branch calculation gives
$x=a+X_1n+O(n^2)$ and $y=a+Y_1n+O(n^2)$ with $Y_1<0$.
Thus $T=-2\lambda Y_1n+O(n^2)>0$ and $y<a$ sufficiently far back.
Let $\alpha=\lambda\Phi(\rho a)/2$. Subtracting the first-order branch
relations gives
\[
 (\lambda+\alpha)(X_1-Y_1)
 =\alpha(a/\lambda-a)-\lambda Y_1>0.
\]
Thus $z>0$ initially on the positive branch. All other inequalities are
strict at the resident end by Lemma~\ref{inc:AN06:v1:residentbounds},
including $v=a/\lambda<(2/5)/(169/400)<11/10$.

\emph{The output-descent differential inequality.}
The exact lag equation on this orbit is
\[
 L'=-\tfrac12(1+nQ^2)L+\tfrac12(1-n)A
 +\tfrac12n(1-n)[\lambda(K(u)+\phi(u))+Q\phi(u)]
 +\tfrac12n^2g(u)Q^2.
\]
Since $A'=\lambda P x'$, direct differentiation of
\eqref{inc:AN06:v1:TZ} gives
\begin{align}
 T'={}&-\tfrac12(1+nQ^2)T\nonumber\\
 &+n(1-n)\left[A(\lambda-Q^2/2)
       +\tfrac\lambda2(K(u)+\phi(u))+\tfrac12Q\phi(u)\right]
       +\tfrac12n^2g(u)Q^2
       -\lambda(1-n)P x'.
 \label{inc:AN06:v1:Tequation}
\end{align}
Assume temporarily the closure of \eqref{inc:AN06:v1:region}, with
$0<n\le1/2$. Then
\[
 P\le\tfrac12+\tfrac25\tfrac7{10}\tfrac34=\frac{71}{100},\quad
 \frac12\le Q\le\frac9{10},\quad v-y\le\frac85.
\]
Also $K(u)+\phi(u)\ge2\phi(0)>3/4$, because its derivative is
$\Phi(u)-u\phi(u)>0$ on $[0,1]$. Since $z\ge0$,
\[
 x'=\tfrac\lambda2 P[n(v-y)-z]
 \le\tfrac\lambda2 nP(v-y).
\]
The coefficient $\lambda-Q^2/2$ is at least $7/400>0$.
Dropping its positive contribution and the positive $g$ term, the remaining
forcing in \eqref{inc:AN06:v1:Tequation}, divided by $n(1-n)$, is at least
\begin{align*}
 c_*&=\frac{169}{400}\frac38+\frac1{20}
       -\frac12\left(\frac{49}{100}\right)^2
                \left(\frac{71}{100}\right)^2\frac85\\
 &=\frac{27902493}{250000000}>\frac1{10}.
\end{align*}
Here $Q\phi(u)/2\ge1/20$ follows from $Q\ge1/2$ and
$\phi(u)\ge\phi(1)>1/5$.
Consequently, throughout this temporary region,
\begin{equation}\label{inc:AN06:v1:Tinequality}
 T'\ge-\tfrac34 T+\tfrac1{10}n(1-n),
\end{equation}
where $T\ge0$ and $(1+nQ^2)/2\le3/4$ were used only to simplify the
negative coefficient. In particular, at $T=0$ the derivative is strictly
positive.

\emph{Simultaneous first-exit closure.}
The other boundary checks are as follows; $y\le a$ follows by integrating
$y'=-T/2\le0$ and the resident limit, rather than being a separate assumed
boundary sign.
At $y=-n$,
\[
 y'+n'\ge n\left[-\frac{\phi(0)}2+\lambda(1-n)\right]>0,
\]
using $K(u)\le1+\phi(0)$ and $\lambda>\phi(0)$.
At $x=3/4$, its bracket is at most
$n(11/10)+(1-n)a-3/4<0$ for $n\le1/2$.
At $v=11/10$, the bracket in $v'$ is negative because
\[
 \frac{K(3\rho/4)}\rho
 \le\frac{8}{13}+\frac38+\frac{63}{800}<\frac{11}{10}.
\]
At $u=1$, $nu+y\ge0$ and $\phi(1)<\lambda$ imply $u'<0$.
At $u=0$,
\[
 u'=\tfrac12[(1-n)\phi(0)-y/2]\ge[\phi(0)-a]/4>0.
\]
At $z=0$, the already established invariant $v>y$ gives
$x'=\lambda n(v-y)P/2>0$, whereas $y'=-T/2\le0$; hence $z'>0$.
Positivity of $v$ is an imported invariant of the selected branch.
The $T$ boundary was checked above. These strictly inward signs exclude
any first exit before or at $n=1/2$, starting sufficiently far back on the
branch. There is no assumption of monotonicity of $x$ or $u$.

\emph{Quantitative integration.}
Integrating \eqref{inc:AN06:v1:Tinequality} from an earlier time and taking
that time to $-\infty$ gives
\[
 T(s)\ge\frac1{10}\int_{-\infty}^s
 e^{-3(s-t)/4}n(t)(1-n(t))\,dt.
\]
The initial term is nonnegative and tends to zero, since $T=O(n)$ at the
resident end. The identities $n'/n=\lambda(1-n)\le\lambda$ and the
monotonicity of $n$ imply, for $t\le s\le0$,
\[
 n(t)(1-n(t))\ge n(s)(1-n(s))e^{-\lambda(s-t)}.
\]
Thus
\[
 T(s)\ge\frac{n(s)(1-n(s))}{10(3/4+\lambda)}
 \ge\frac5{62}n(s)(1-n(s)).
\]
At $s=0$, $n=1/2$ and $y'=-T/2$, proving
\eqref{inc:AN06:v1:Tlower}.
\end{proof}

\section{A negative total rate at half-learning}
\label{inc:AN06:v1:halfdescent}
\begin{theorem}[Uniform half-mass descent witness]
\label{inc:AN06:v1:halfdescenttheorem}
For every $\rho\in\mathcal K$ and $\zeta\in[4,5]$, the selected orbit is
probe-active at $s=0$ and
\begin{equation}\label{inc:AN06:v1:halfrate}
 \mathcal E_\zeta'(0)\le-\frac{152833}{31000000}<-\frac1{250}.
\end{equation}
Also $\mathcal E_\zeta(0)<0$.
\end{theorem}
\begin{proof}
At $s=0$, write $B=v-y\le8/5$. Completing the square gives
\[
 (x-y)x'=\frac\lambda2 Pz(nB-z)
 \le\frac{\lambda Pn^2B^2}{8}
 \le\frac{3479}{125000}.
\]
This does not require $x'\ge0$. By Theorem~\ref{inc:AN06:v1:regiontheorem},
$\zeta-x\ge13/4$ and $y'\le-5/496$. Therefore
\[
 \mathcal E_\zeta'(0)
 \le\frac{3479}{125000}-\frac{13}{4}\frac5{496}
 =-\frac{152833}{31000000}<-\frac1{250}.
\]
Finally,
$\mathcal E_\zeta=(\zeta-x)[y-(\zeta+x)/2]$.
Here $x>y>-1/2$, $y<2/5$, and $\zeta\ge4$, so both factors have the
signs needed for strict negativity.
\end{proof}

\begin{corollary}[A crossover strictly after half-mass]
\label{inc:AN06:v1:posthalf}
For each $\rho\in\mathcal K$, $\zeta\in[4,5]$, there is a finite
negative-to-positive zero of $\mathcal E_\zeta'$ at a time $s_*>0$ before
the first strong gate hit. A sufficiently small closed interval about this
zero has $n>1/2$, $d_1<0<d_2$ everywhere, and ordered positive-width
subintervals of strictly negative and strictly positive total rate.
\end{corollary}
\begin{proof}
The imported gate-hit theorem gives a finite first hit $t_g>0$ of $x=\zeta$;
the half-region bounds exclude such a hit before $s=0$.
At $s=0$ the error is negative and its derivative is negative, while its
value at the hit is zero. There is a point of positive derivative in
$(0,t_g)$. On this interval the derivative of the active error is real
analytic. Its first transition to positivity therefore occurs at an
isolated zero $s_*>0$, with a negative sign just to the left and positive
sign just to the right. This zero has finite odd order, but need not be
simple. At the zero, $d_1=-(\zeta-x)L/2<0$, so $d_2=-d_1>0$.
Continuity supplies a common source-sign neighborhood, and analyticity
supplies the two total-rate windows. Since the entire neighborhood can be
chosen inside $(0,t_g)$, $n>1/2$ on it.
\end{proof}

\section{A gate-aware population rate estimate}
\label{inc:AN06:v1:populationrobustness}
This section supplies the quantitative stability interface used for common
windows. It does not assume that an initialized population has reached it.
Let $\mu_s,\mu_w$ be nonnegative finite measures in the scaled charts and put
\begin{gather*}
 M=\mu_s(\R^2),\quad N=\mu_w(\R^2),\quad
 Y=\int y\,d\mu_s,\quad V=\int v\,d\mu_w,\\
 K_s=\int K(\rho x)\,d\mu_s,\qquad K_w=\int K(u)\,d\mu_w.
\end{gather*}
The overlap kernel is
\[
 F(t,r)=\phi(\min(t,r))+r\Phi(\min(t,r)),\quad
 \mathcal S(x)=\int F(\rho x,\rho\widetilde x)\,d\mu_s,
 \quad \mathcal W(u)=\int F(u,\widetilde u)\,d\mu_w.
\]
The strong source velocities and weak velocities are
\begin{align}
 x'_1&=-\tfrac12(1-M)x,&
 y'_1&=-\tfrac12[(1-M)y+Y+K_w],
 \label{inc:AN06:v1:popsystem}\\
 x'_2&=\tfrac12[\rho\phi(\rho x)-\rho\mathcal S(x)
       +\lambda\{V+(1-N)y\}\Phi(\rho x)],&
 y'_2&=\tfrac\rho2(1-N)K(\rho x),\nonumber\\
 u'&=\tfrac12[\phi(u)-\mathcal W(u)-\{Y+(1-M)v\}\Phi(u)
                  -\lambda(1-N)u],\nonumber\\
 v'&=\tfrac12[\rho K_s-\lambda V-\lambda(1-N)v-(1-M)K(u)].\nonumber
\end{align}
Strong reaction is $1-M$ and weak reaction is $\lambda(1-N)$.
Thus $M'=M(1-M)$ and $N'=\lambda N(1-N)$.
Use
\[
 A_\zeta=\int(\zeta-x)_+\,d\mu_s,\qquad
 B_\zeta=\int y(\zeta-x)_+\,d\mu_s,\qquad
 \mathcal E=\tfrac12A_\zeta^2-\zeta A_\zeta+B_\zeta.
\]
The use of $A_\zeta$ here is distinct from $A=\rho K(\rho x)$ in Section 2.

\begin{lemma}[Rates including nonunit strong mass]
\label{inc:AN06:v1:generalrates}
At points/times where the characteristic chain rule applies, with the
indicator interpreted as $\1_{x<\zeta}$, the leading source rates satisfy
\begin{align}
 d_1={}&-\tfrac12A_\zeta(Y+K_w)+(1-M)\mathcal C,\nonumber\\
 \mathcal C={}&(A_\zeta-\zeta)A_\zeta+\tfrac12B_\zeta
 -\tfrac12\int_{x<\zeta}(\zeta-A_\zeta-y)x\,d\mu_s,
 \label{inc:AN06:v1:generalrates-eq}\\
 d_2={}&\int_{x<\zeta}(\zeta-A_\zeta-y)x'_2\,d\mu_s
       +\int(\zeta-x)_+y'_2\,d\mu_s.
 \nonumber
\end{align}
Their sum is $\mathcal E'$ almost everywhere along compact, regular
characteristic solutions.
\end{lemma}
\begin{proof}
For source 1,
\[
 (A_\zeta)'_1=(1-M)A_\zeta-\int_{x<\zeta}x'_1\,d\mu_s,
\]
\[
 (B_\zeta)'_1=(1-M)B_\zeta+
 \int_{x<\zeta}[(\zeta-x)y'_1-yx'_1]\,d\mu_s.
\]
Substitute these expressions into
$d_1=(A_\zeta-\zeta)(A_\zeta)'_1+(B_\zeta)'_1$ and use
\eqref{inc:AN06:v1:popsystem}. Source 2 has no strong reaction term and
gives the displayed $d_2$ directly.
For the almost-everywhere total identity, apply the positive-part chain
rule to each characteristic and then Fubini; compact support bounds the
integrands. At a gate the chosen selector is zero. The set of times at
which a given absolutely continuous characteristic equals the gate and
has nonzero total normal velocity has measure zero. No assertion that
there is no boundary mass at every time is required. These are the
selector-defined leading source rates, not a claimed pointwise
finite-noise source-transfer theorem at boundary atoms.
\end{proof}

Let a comparison state be the coherent measure
$\delta_{(X,Y_*)}$ and $n_*\delta_{(U,V_*)}$, where $0\le n_*\le1$.
Define the mass-weighted first-moment distance
\begin{align}
 \mathfrak d={}&|M-1|+|N-n_*|\nonumber\\
 &+\int(|x-X|+|y-Y_*|)\,d\mu_s
 +\int(|u-U|+|v-V_*|)\,d\mu_w.
 \label{inc:AN06:v1:distance}
\end{align}
The coordinates $Y_*,V_*$ here are atom coordinates, not total moments.

\begin{proposition}[Explicit rate robustness with gate mismatches]
\label{inc:AN06:v1:ratebound}
Suppose $R\ge2$, $0<\rho\le1$, $M,N\le2$, both population supports and
the comparison coordinates lie in $[-R,R]^2$, $|\zeta|\le R$, and
$\zeta-X\ge g_0>0$. At the same $\rho,\zeta$,
\begin{equation}\label{inc:AN06:v1:ratebound-eq}
 |d_p-d_{p,*}|\le C(R,g_0)\mathfrak d\quad(p=1,2),\qquad
 C(R,g_0)=64R^2(1+g_0^{-1}).
\end{equation}
Also $|\mathcal E-\mathcal E_*|\le10R^2\mathfrak d$.
The population need not have every strong label probe-active.
\end{proposition}
\begin{proof}
Write $D_s,D_w$ for the two integral terms in
\eqref{inc:AN06:v1:distance}. The elementary bounds
$K\le2R$, $|\phi'|\le1$, $K'\le1$, and
$|\partial_tF|\le2R$, $|\partial_rF|\le1$ hold on the relevant arguments.
They imply
\begin{gather*}
 |A_\zeta-(\zeta-X)|\le2R\mathfrak d,\quad
 |B_\zeta-Y_*(\zeta-X)|\le2R^2\mathfrak d,\\
 |Y-Y_*|\le R\mathfrak d,\quad
 |K_w-n_*K(U)|\le2R\mathfrak d,\quad
 |V-n_*V_*|\le R\mathfrak d.
\end{gather*}
For a strong receiver, writing $D=|x-X|+|y-Y_*|$, the donor/receiver split
of $\mathcal S$ gives
\[
 |x'_2-x'_{2,*}|\le2R(D+\mathfrak d),\qquad
 |y'_2-y'_{2,*}|\le\tfrac12D+R\mathfrak d.
\]
Indeed the donor and mass changes in $\mathcal S$ cost at most
$2R\mathfrak d$, and its reference-receiver change costs at most
$2R|x-X|$. The other terms follow from the displayed moment bounds and
$|n_*V_*+(1-n_*)Y_*|\le R$.
Uniformly $A_\zeta\le4R$, $|Y+K_w|\le6R$, $|x'_2|\le4R$, and
$|y'_2|\le R$.

For $d_1$, the lag product costs at most $9R^2\mathfrak d$.
In \eqref{inc:AN06:v1:generalrates-eq},
$|\mathcal C|\le20R^2+2R^2+6R^2=28R^2$.
Thus $|d_1-d_{1,*}|\le37R^2\mathfrak d$.

For $d_2$, put $c=\zeta-A_\zeta-y$ and $c_*=X-Y_*$.
Integrating the smooth difference $cx'_2-c_*x'_{2,*}$ over active
receivers costs at most $30R^2\mathfrak d$.
The remaining active-mass difference is bounded using
\[
 \mu_s\{x\ge\zeta\}\le D_s/g_0,\qquad
 |\mu_s\{x<\zeta\}-1|\le|M-1|+D_s/g_0.
\]
Since $|c_*x'_{2,*}|\le8R^2$, this costs at most
$8R^2(1+g_0^{-1})\mathfrak d$.
The output-rotation term costs at most $7R^2\mathfrak d$:
its receiver difference is bounded by $2RD+2R^2\mathfrak d$, and its
reference mass term by $2R^2|M-1|$.
The total is at most $(45+8/g_0)R^2\mathfrak d$, below the stated constant.
Finally the error difference is bounded by
$4R|A_\zeta-(\zeta-X)|+|B_\zeta-Y_*(\zeta-X)|$.
All integrals have compact support; the inactive-mass estimate includes
labels exactly on the gate.
\end{proof}

\begin{lemma}[Coherent parameter sensitivity at a fixed state]
\label{inc:AN06:v1:paramrate}
For a fixed coherent state satisfying the bounds of
Proposition~\ref{inc:AN06:v1:ratebound}, while its strong gate stays active,
changing only the parameters gives
\[
 |\Delta d_p|\le8R^3|\Delta\rho|+2R|\Delta\zeta|\quad(p=1,2).
\]
\end{lemma}
\begin{proof}
The $d_1$ formula is independent of $\rho$ and has $\zeta$ derivative
$-L/2$, of absolute value at most $3R/2$.
In $d_2$ the input-rotation term is independent of $\zeta$ and the
output-rotation term has $\zeta$ derivative of absolute value at most $R$.
For the $\rho$ derivative write
$x'_2=\rho^2\Phi(\rho X)[n_*V_*-X+(1-n_*)Y_*]/2$.
Its derivative has absolute value at most $2R^2$; multiplying by
$|X-Y_*|\le2R$ costs $4R^3$.
The derivative of the output-rotation term costs at most $3R^2$.
Their sum is below $8R^3$. Integrate along the parameter segment.
\end{proof}

\section{Common time windows and a nonempty parameter rectangle}
\label{inc:AN06:v1:commonwindows}
We now define every constant by an analytical operation on the selected
orbit. No decimal approximation to a crossing, gate time, or minimum is
used. The resulting widths are positive but are not evaluated numerically
in this increment.

\begin{lemma}[Smooth parameter dependence of the selected branch]
\label{inc:AN06:v1:branchsmooth}
The map $(\rho,s)\mapsto Z_\rho(s)$ is continuously differentiable in
$\rho$ on $\mathcal K$ and every compact finite time interval. Its
$\rho$ derivative and the state itself are bounded there.
\end{lemma}
\begin{proof}
The implicit resident roots are smooth since their defining derivatives
are nonzero. The four angular eigenvalues at the resident are strictly
negative; on the compact set $\mathcal K$ they are bounded away from zero.
The unstable mass eigenvalue is $\lambda\ge169/400$.
Here is a local parameter argument which also fixes the time phase.
Let $w$ be the angular displacement from the resident, use $n$ as the
independent variable, and write
\[
 n\frac{dw}{dn}=B_\rho w+H_\rho(n,w),\qquad
 B_\rho=J_{\mathrm{ang},\rho}/\lambda,
\]
where $H_\rho(n,w)=O(n)+O(n|w|+|w|^2)$ locally, uniformly in $\rho$.
The negative spectrum gives constants $C,\kappa>0$ with
$\|(n/t)^{B_\rho}\|\le C(t/n)^\kappa$ for $0<t\le n$.
The incoming branch solves
\[
 w(n)=\int_0^n(n/t)^{B_\rho}H_\rho(t,w(t))\,\frac{dt}{t}.
\]
In a ball $\sup_{0<n\le n_0}|w(n)|/n\le K$, this map preserves the ball
for sufficiently large fixed $K$ and sufficiently small $n_0$, and its
Lipschitz constant is $O(n_0)$. This follows directly by integrating
$(t/n)^\kappa t$ and $(t/n)^\kappa t^2$ against $dt/t$.
It is therefore a uniform contraction.
Differentiating the matrix kernel in $\rho$ introduces at most an
additional constant times $1+|\log(n/t)|$, with a possibly smaller
positive exponent $\kappa$. This is integrable against the same powers
of $t$. Differentiation of the contraction equation therefore yields a
unique continuous parameter derivative, by the inverse of $I-D_w\mathcal T$.
The resulting solution is the unique incoming branch of
\texttt{pa:branch}. Starting at $n=n_0$, the ordinary smooth ODE and its
parameter variational equation continue this dependence to every finite
$s$, with phase fixed by $n(s)=(1+e^{-\rho^2s})^{-1}$.
Global forward continuation of the branch is an imported result; local
ODE continuation and compactness in $\rho$ give the asserted finite-window
bounds.
\end{proof}

\paragraph{Central orbit and an analytic zero.}
Set $\rho_c=2/3$, $\zeta_c=9/2$, and write
$Z_c(s)=Z_{\rho_c}(s)$. Let $t_g$ be its first gate hit after zero.
On $[0,t_g]$, define $h(s)$ by the derivative of the active polynomial
error formula \eqref{inc:AN06:v1:error}; this formula is smooth up to the
endpoint. Define
\[
 s_*:=\inf\{s\in(0,t_g):h(s)>0\}.
\]
Corollary~\ref{inc:AN06:v1:posthalf} proves $0<s_*<t_g$, with an isolated
negative-to-positive zero at $s_*$. Let $k\ge1$ be its finite odd order and
\[
 b=\frac{h^{(k)}(s_*)}{k!}>0,\qquad
 g_*=\zeta_c-x_c(s_*)>0,\qquad
 p_*=\frac12g_*L_c(s_*)>0.
\]
Define the finite derivative bounds
\begin{align*}
 C_g&=1+\sup_{[0,t_g]}|x_c'|,\\
 C_d&=1+\sup_{[0,t_g]}(|d_{1,c}'|+|d_{2,c}'|),\\
 C_h&=1+\frac1{(k+1)!}\sup_{[0,t_g]}|h^{(k+1)}|.
\end{align*}
All source derivatives in these definitions use the active formulas.
Set
\begin{equation}\label{inc:AN06:v1:rdef}
 r=\min\left\{\frac{s_*}{4},\frac{t_g-s_*}{4},\frac14,
 \frac{g_*}{4C_g},\frac{p_*}{4C_d},\frac b{2C_h}\right\}>0,
\end{equation}
\begin{equation}\label{inc:AN06:v1:windows}
 J=[s_*-r,s_*+r],\quad
 J_-=[s_*-r,s_*-r/2],\quad
 J_+=[s_*+r/2,s_*+r],\quad
 \gamma_0=\frac{br^k}{2^{k+1}}>0.
\end{equation}
On $J$, the central gap is at least $3g_*/4$,
$d_{1,c}\le-3p_*/4$, and $d_{2,c}\ge3p_*/4$.
Taylor's formula and \eqref{inc:AN06:v1:rdef} give
\begin{equation}\label{inc:AN06:v1:centralmargins}
 h\le-\gamma_0\text{ on }J_-,\qquad h\ge\gamma_0\text{ on }J_+.
\end{equation}
These inequalities are quantitative even when the zero is not simple:
$|h(s_*+t)-bt^k|\le C_h|t|^{k+1}\le(b/2)|t|^k$.

\paragraph{Explicit local parameter widths.}
Define
\begin{align}
 R&=2+\max\left\{5,\sup_{\rho\in\mathcal K,\,s\in J}
                  \norm{Z_\rho(s)}_\infty\right\},\nonumber\\
 S_\rho&=1+\sup_{\rho\in\mathcal K,\,s\in J}
                  \norm{\partial_\rho Z_\rho(s)}_1,\nonumber\\
 \sigma&=n_c(s_*-r)-\tfrac12>0,\qquad
 C_0=C(R,g_*/2),\qquad \eta=\min\{p_*/8,\gamma_0/8\},
 \label{inc:AN06:v1:constants}\\
 \delta_\rho&=\min\left\{\frac1{120},
   \frac{\eta}{2(C_0S_\rho+8R^3)},
   \frac{\sigma}{2S_\rho},\frac{g_*}{4S_\rho}\right\},\nonumber\\
 \delta_\zeta&=\min\left\{\frac14,\frac{g_*}{4},\frac{\eta}{4R}\right\}.
 \nonumber
\end{align}
Lemma~\ref{inc:AN06:v1:branchsmooth} makes $R,S_\rho$ finite. All the
other constants were proved positive. Thus the closed rectangle
\begin{equation}\label{inc:AN06:v1:rectangle}
 \mathcal P_*=[\rho_c-\delta_\rho,\rho_c+\delta_\rho]
 \times[\zeta_c-\delta_\zeta,\zeta_c+\delta_\zeta]
\end{equation}
is nonempty with nonzero side lengths, lies in $\mathcal K\times[4,5]$,
and contains $(2/3,9/2)$ in its interior.
These are orbit-defined analytical endpoints, not numerical enclosures.

\begin{theorem}[Common learned competition witnesses]
\label{inc:AN06:v1:uniformtheorem}
For every $(\rho,\zeta)\in\mathcal P_*$, the \emph{same} windows
\eqref{inc:AN06:v1:windows} on the half-mass-centered clock satisfy
\begin{align}
 n_\rho&\ge\tfrac12+\tfrac\sigma2,
 &\zeta-x_\rho&\ge\tfrac{g_*}{4} &&\text{on }J,\nonumber\\
 d_1&\le-\tfrac58p_*,& d_2&\ge\tfrac58p_* &&\text{on }J,
 \label{inc:AN06:v1:uniformmargins}\\
 \mathcal E_\zeta'&\le-\tfrac34\gamma_0&&&\text{on }J_-,\nonumber\\
 \mathcal E_\zeta'&\ge\tfrac34\gamma_0&&&\text{on }J_+.\nonumber
\end{align}
\end{theorem}
\begin{proof}
For the two coherent states at a fixed $s\in J$, their distance
\eqref{inc:AN06:v1:distance} is at most $S_\rho|\rho-\rho_c|$,
since both weak masses are at most one.
The central comparison gate at the new $\zeta$ is at least $g_*/2$.
Apply Proposition~\ref{inc:AN06:v1:ratebound} at this new parameter and
then Lemma~\ref{inc:AN06:v1:paramrate} to the central state. The difference
of each source rate from its central value is at most
\[
 (C_0S_\rho+8R^3)|\rho-\rho_c|+2R|\zeta-\zeta_c|\le\eta.
\]
The source margins and the total-rate margins follow from
\eqref{inc:AN06:v1:centralmargins}, using
$\eta\le p_*/8$ and $2\eta\le\gamma_0/4$.
The mass and gate bounds follow directly from the last two restrictions
in $\delta_\rho$ and the restriction $\delta_\zeta\le g_*/4$.
\end{proof}

\section{The population tolerance required on these windows}
\label{inc:AN06:v1:tolerance}
Define
\begin{equation}\label{inc:AN06:v1:epsdef}
 C_1=C(R,g_*/4),\qquad
 \varepsilon_*=
 \min\left\{\frac14,\frac\sigma4,\frac{p_*}{8C_1},
                         \frac{\gamma_0}{8C_1}\right\}>0.
\end{equation}

\begin{theorem}[Conditional population witness transfer]
\label{inc:AN06:v1:populationtransfer}
Fix $(\rho,\zeta)\in\mathcal P_*$. Suppose a solution of the defined
leading population system has $M,N\le2$ and chart supports in $[-R,R]^2$
throughout $J$, and its distance to the selected coherent state at the
same $(\rho,s)$ satisfies
\begin{equation}\label{inc:AN06:v1:tolerancecondition}
 \sup_{s\in J}\mathfrak d(s)\le\varepsilon_*.
\end{equation}
Then, with the selector convention of
Lemma~\ref{inc:AN06:v1:generalrates},
\begin{align}
 M&\ge\tfrac34,\qquad N\ge\tfrac12+\tfrac\sigma4,
 &M,N&\le\tfrac54 &&\text{on }J,\nonumber\\
 d_1&\le-\tfrac12p_*,&d_2&\ge\tfrac12p_* &&\text{on }J,
 \label{inc:AN06:v1:populationmargins}\\
 d_1+d_2&\le-\tfrac12\gamma_0&&&\text{on }J_-,\nonumber\\
 d_1+d_2&\ge\tfrac12\gamma_0&&&\text{on }J_+.\nonumber
\end{align}
For a regular characteristic solution, the last two inequalities hold
for the derivative of $\mathcal E$ almost everywhere. There is an
interior minimizer $s_{\min}\in(s_*-r,s_*+r)$ such that
\begin{equation}\label{inc:AN06:v1:rebound}
 \mathcal E(s_*-r)-\mathcal E(s_{\min})\ge\frac{\gamma_0r}{4},\qquad
 \mathcal E(s_*+r)-\mathcal E(s_{\min})\ge\frac{\gamma_0r}{4}.
\end{equation}
\end{theorem}
\begin{proof}
The mass estimates follow from Theorem~\ref{inc:AN06:v1:uniformtheorem}
and the mass terms in $\mathfrak d$.
Apply Proposition~\ref{inc:AN06:v1:ratebound} with reference gap $g_*/4$.
Each source-rate change is at most $C_1\varepsilon_*$, hence at most
$p_*/8$ and at most $\gamma_0/8$. Subtract these budgets from
\eqref{inc:AN06:v1:uniformmargins}. Finally each signed subinterval has
length $r/2$. Integrating proves the two inequalities, and excludes the
endpoints as minimizers, exactly as in the checkpoint's
\texttt{pa:assembly}.
\end{proof}

\begin{remark}[What the metric controls and what it does not]
The tolerance includes both masses and first absolute moments in
\emph{both} angular coordinates of both populations. It is not a
strong-family variance condition alone. A strong input variance bound can
supply part of it by Cauchy--Schwarz, but the means, weak coordinates,
and masses must also be controlled. Gate mismatch costs $D_s/g_0$ and is
not discarded. The compact support condition is substantive; population
outside these charts is not covered by this theorem. Separate bounds on
that population's output and rate contributions are still required for
an initialized application.
\end{remark}

\section{A finite-window propagation estimate, not resident passage}
\label{inc:AN06:v1:finitewindow}
\begin{proposition}[Conditional propagation in the defined leading system]
\label{inc:AN06:v1:propagation}
Let a leading population solution and the selected coherent solution at
the same $\rho$ remain in the compact class used in
Proposition~\ref{inc:AN06:v1:ratebound} on $[a,b]$.
Then their distance \eqref{inc:AN06:v1:distance} satisfies
\begin{equation}\label{inc:AN06:v1:propagationbound}
 \mathfrak d(s)\le e^{64R(s-a)}\mathfrak d(a)
 \quad(a\le s\le b).
\end{equation}
Consequently, conditional on chart validity on $J$, an entry tolerance
$\mathfrak d(s_*-r)\le\varepsilon_*e^{-128Rr}$ suffices for
\eqref{inc:AN06:v1:tolerancecondition}.
\end{proposition}
\begin{proof}
For masses in $[0,2]$, the difference of the two logistic mass fields is
bounded by three times the mass difference. Relative weights of fixed
labels within either family are constant in the leading system.
The receiver/donor estimates in the proof of
Proposition~\ref{inc:AN06:v1:ratebound}, also applied to $u',v'$, give
\[
 |x'-X'|+|y'-Y_*'|\le8R(|x-X|+|y-Y_*|+\mathfrak d),
\]
\[
 |u'-U'|+|v'-V_*'|\le8R(|u-U|+|v-V_*|+\mathfrak d).
\]
For clarity, the weak overlap difference costs at most
$2R\mathfrak d+2R|u-U|$; the remaining weak terms are controlled by
$|Y-Y_*|$, $|V-n_*V_*|$, $|M-1|$, and $|N-n_*|$.
The strong and weak reaction coefficients have absolute value at most
one. Differentiating the two weighted distance integrals in the
almost-everywhere/Dini sense therefore gives
\[
 D^+D_s\le(1+8R)D_s+16R\mathfrak d,\qquad
 D^+D_w\le(1+8R)D_w+16R\mathfrak d.
\]
Adding the mass bounds yields
$D^+\mathfrak d\le(3+40R)\mathfrak d\le64R\mathfrak d$.
Integration proves the claim. Compactness justifies integration of the
characteristic inequalities and also covers zero family mass by
continuity. This argument presupposes, and does not prove, chart validity.
\end{proof}

The exponent in \eqref{inc:AN06:v1:propagationbound} is deliberately
coarse and is used only on the fixed witness interval of length $2r$.
It is \emph{not} a substitute for AN3's competing-rate estimate over a
delay of order $\log(1/S_0)$, nor does it close strong splitting or entry.

\section{Static exact-Gaussian learning bounds}
\label{inc:AN06:v1:staticlearning}
Let the entire state be represented by the same two measures, now embedded
in the exact finite-noise network as
\begin{align*}
 a_s&=(\cos(qx),\sin(qx)),& b_s&=(\cos(qy),\sin(qy)),\\
 a_w&=(\sin(qu),\cos(qu)),& b_w&=(-\sin(qv),\cos(qv)),\\
 f(X)&=\int b_s(a_s^\top X)_+\,d\mu_s
       +\int b_w(a_w^\top X)_+\,d\mu_w.
\end{align*}
This is an exact state representation, not an evolution claim.
Use $P_1=\mathcal N(e_1,q^2I_2)$ and
$P_2=\mathcal N(\rho e_2,q^2I_2)$,
$Q_1=1+q^2$, and $Q_2=\rho^2+q^2$.

\begin{proposition}[Responses and losses from a fully charted state]
\label{inc:AN06:v1:staticproposition}
Suppose $\rho\in\mathcal K$, $0<q\le1$, both supports lie in $[-R,R]^2$,
$R\ge2$, and $M,N\le2$. Define the exact observables
\[
 \mathsf m_p=\frac{\E_{P_p}[X_pf_p(X)]}{Q_p},\qquad
 \mathcal L_p=\tfrac12\E_{P_p}\|f(X)-X\|^2.
\]
Writing $M_1=M$, $M_2=N$, one has
\begin{equation}\label{inc:AN06:v1:staticbounds}
 |\mathsf m_p-M_p|\le25qR,\qquad
 \sqrt{2\mathcal L_p}\le |1-M_p|\sqrt{Q_p}+17qR.
\end{equation}
If the mass conclusions of
Theorem~\ref{inc:AN06:v1:populationtransfer} hold, $0<S_0\le1$, and
\begin{equation}\label{inc:AN06:v1:qlearn}
 q\le q_{\rm learn}:=
 \min\left\{1,\frac{\sigma}{100R},\frac1{100R},
                  \frac{13/20}{136R}\right\},
\end{equation}
then
\begin{equation}\label{inc:AN06:v1:learningconclusion}
 \mathsf m_1,\mathsf m_2\ge\frac12,\qquad
 \mathcal L_p\le\frac{25}{36}\mathcal L_p(0),
\end{equation}
where $\mathcal L_p(0)$ is the loss of the specified isotropic initial
output $f_0(X)=S_0X/4$ on the same Gaussian cluster.
\end{proposition}
\begin{proof}
On cluster 1, a strong feature obeys
\[
 \|b_s(a_s^\top X)_+-e_1X_1\|_{L^2(P_1)}
 \le2qR\|X\|_{L^2(P_1)}+\|(-X_1)_+\|_{L^2(P_1)}\le5qR.
\]
This uses the Lipschitz property of the positive part,
$\|a_s-e_1\|,\|b_s-e_1\|\le qR$,
$\|X\|_{L^2(P_1)}<2$, and
$(-1-qZ)_+\le q|Z|$.
A weak feature on cluster 1 has norm at most
\[
 |\sin(qu)|\|X_1\|_{L^2}+|\cos(qu)|\|X_2\|_{L^2}\le3qR.
\]
Minkowski and $M,N\le2$ yield
\[
 \|f-e_1MX_1\|_{L^2(P_1)}\le16qR.
\]
The identical calculation with the roles of the axes interchanged gives
$\|f-e_2NX_2\|_{L^2(P_2)}\le16qR$; here
$(-\rho-qZ)_+\le q|Z|$.
Cauchy--Schwarz and $\sqrt{Q_p}\ge13/20$ give the response bound,
since $16/(13/20)<25$. Adding the unused target coordinate, whose
$L^2$ norm is $q$, gives the loss bound.

For the last assertion, $M\ge3/4$ and $N\ge1/2+\sigma/4$.
The restrictions $25qR\le1/4$ and $25qR\le\sigma/4$ prove the response
thresholds. Both masses lie between $1/2$ and $5/4$, so
$|1-M_p|\le1/2$. Also $17qR/\sqrt{Q_p}\le1/8$, hence
$\sqrt{2\mathcal L_p}\le(5/8)\sqrt{Q_p}$.
The actual initial cluster residual has norm
\[
 \sqrt{2\mathcal L_p(0)}=(1-S_0/4)\sqrt{Q_p+q^2}
 \ge(3/4)\sqrt{Q_p}.
\]
Taking the squared ratio gives $25/36$.
Gaussian negative tails were bounded, not declared zero.
\end{proof}

This proposition does not require the exact finite-noise flow to follow
the leading flow. It only says that a state meeting the displayed
hypotheses is genuinely learned in the required observables. Proving that
the initialized trajectory reaches and remains in that state class, and
that its exact source rates meet the leading signed budgets, remains open.

\section{Open residue and implications for the next increment}
\label{inc:AN06:v1:residue}
\paragraph{What is closed in the limiting witness problem.}
There is now a rational, parameter-uniform descent witness at weak
half-mass, not merely a negative slope as $n\downarrow0$.
The subsequent crossing is strictly after half-mass. Source signs hold
on the whole common interval $J$, not only at the zero, and the windows,
parameter widths, population tolerance, and rebound lower bound are
specified by positive analytical constants. No crossing uniqueness or
transversality is assumed.
The common-window widths remain unevaluated analytic orbit functionals;
they are not certified numerical values or a claim of a large rectangle.

\paragraph{What remains open.}
The initialized argument must still prove coupled strong contraction and
weak capture; quantify the weak seed and the full strong shape;
control all population outside the charts; perform nonlinear resident
passage; and transfer the original finite-noise rates with errors smaller
than the window margins. The bounded first-moment tolerance is an explicit
\emph{target}, not an initialized conclusion. The finite-window estimate
does not handle the diverging entry delay. The static learning estimate
does not supply a finite-noise dynamical approximation.
The range $\zeta\in[4,5]$ was chosen to make the half-learning inequality
close with elementary constants; descent at $\zeta\simeq1.745$ and
canonical positive-noise coverage are not established here.

\paragraph{A joint-scaling correction for the arrival roadmap.}
The condition $q^2\log(1/S_0)\to0$ does not by itself select the
third-quadrant reservoir regime described in the user's review.
For example $S_0=q^3$ satisfies this condition and even $S_0/q^2\to0$,
but for $\rho=2/3$ it has
\[
 \frac{\log(\delta/S_0)}{\log(1/q)}\longrightarrow3<\frac2\lambda=\frac92.
\]
It is therefore on the bulk side of the review's proposed arrival
crossover scaling. The crossover description itself still needs the
corresponding uniform transport estimates. A sufficient scale separation
\emph{for that proposed reservoir timing criterion} is
\[
 \frac{\log(1/S_0)}{\log(1/q)}\to\infty,\qquad
 q^2\log(1/S_0)\to0;
\]
$S_0=e^{-1/q}$ satisfies both. These scale relations do not prove capture.
The Riccati refinement of AN02 and its exponentially small exponent loss
remain a proposed refinement, not a result of this file.

\paragraph{Proposed integration location after approval.}
Insert Sections 2--3 after \texttt{pa:limitreversal}, updating only its
learning-witness limitation for the proved offset range.
Sections 4--7 supply a quantitative conditional interface for AN3/AN6/AN8.
Section 8 belongs with AN5's mass-to-exact-observable interface.
Do not replace AN2/AN3/AN7's open status or merge this file without review.

\end{document}
